\chapter{Related work and research context}
\label{ch:related}

This chapter positions the investigation within four lines of existing work,
not as a survey but to establish which parts of the problem are recognised
open questions elsewhere and which parts this report investigates directly.

\textbf{Parametric articulated models.} \term{SMPL}~\cite{smpl} learns a
human shape and pose space from a large registered scan corpus under a fixed
body topology. \term{SMAL}~\cite{smal}, from the same \emph{3D Menagerie}
line of work, relaxes the cooperative-subject and topology assumptions for
quadrupeds by learning a shape space from toy figurines and fitting it to
images and video~\cite{smalify}. \term{smilify} continues this progression
to insects, where the body plan itself differs and every structure of
interest is thin and articulated (Chapter~\ref{ch:pipeline}).

\textbf{Point-cloud and mesh registration.} Classical iterative closest
point and its variants align point sets by alternating nearest-neighbour
correspondence and rigid transformation estimation, and are known to
converge to spurious local optima under poor initialisation or partial
overlap. Learning-based registration addresses this directly: \term{Deep
Closest Point}~\cite{dcp} learns point embeddings for closest-point
matching end-to-end; \term{PREDATOR}~\cite{predator} predicts an overlap
region between two point clouds before matching, specifically to handle
low-overlap pairs; and \term{GeoTransformer}~\cite{geotransformer} combines
geometric attention with coarse-to-fine matching for dense correspondence.
These methods establish that initialisation, overlap, and correspondence
are recognised as separable, fundamental difficulties in point-cloud
registration generally. This report studies a harder, model-based variant:
registration constrained to the image of a parametric articulated model
rather than a free rigid or non-rigid point-set alignment, where the
nearest-surface matching implicit in a Chamfer-style objective is the local
matching mechanism and the open question is what information that mechanism
lacks (\S\ref{sec:registration}).

\textbf{Shape correspondence.} Beyond point-to-point matching,
\term{functional maps}~\cite{fmaps} represent a correspondence between two
shapes as a compact linear map between functions defined on them, and a
substantial subsequent literature extends this to partial and non-isometric
settings; \term{optimal transport}~\cite{ot} provides a complementary
globally-consistent matching formulation under an explicit transport cost,
with partial-transport variants suited to incomplete observations. These
are the two directions this report's outlook (\S\ref{sec:learned-correspondence},
Chapter~\ref{ch:conclusion}) points to for closing the local-to-global gap
exposed by the learned-descriptor experiments.

\textbf{Defective mesh repair.} \term{Alpha Wrap}~\cite{alphawrap}
constructs a watertight offset surface enclosing arbitrarily defective
input, with an explicit offset parameter trading fidelity against
robustness; \term{ManifoldPlus}~\cite{manifoldplus} instead repairs a given
triangle soup directly into a watertight manifold. Both target exactly the
defect classes (holes, self-intersections, non-manifold geometry)
present in the scan corpus this report works with (\S\ref{sec:preprocessing}).

None of these methods is claimed to solve the problem studied here; they are
cited to establish that each mechanism this report investigates
(initialisation, correspondence, overlap-like partiality, and defective
input) is an independently recognised difficulty in 3D registration, and
that this report's contribution is to characterise how these difficulties
interact specifically for an articulated, model-based, biological fitting
problem.
